% ledger: a ruled record sheet for exercises and exams. Libertinus type when
% installed, compact bold headings, a ruled header block with write-in fields
% in place of a title page, small running heads between hairlines, soft grey
% callouts, and a "Continued next page" footer. Designed for scrartcl on a
% 10pt page with 18mm margins.
%
% Cover fields: title and label share the first row, series and edition the
% second, subtitle and author the third. `code` leads the running head and
% `fields` lists write-in labels (for example Student ID: and Name:).
\definecolor{pressink}{HTML}{000000}
\definecolor{pressstrong}{HTML}{000000}
\definecolor{pressaccent}{HTML}{929292}
\definecolor{pressaccentdark}{HTML}{5C5C5C}
\definecolor{pressaccentlight}{HTML}{C8C8C8}
\definecolor{presssecondary}{HTML}{0000FF}
\definecolor{pressmuted}{HTML}{5C5C5C}
\definecolor{pressrule}{HTML}{000000}
\definecolor{presssoft}{HTML}{F6F6F6}

\makeatletter
% Libertinus replaces the format's fonts only where it is installed, so a
% machine without it still renders with the fonts set in YAML.
\@ifpackageloaded{fontspec}{%
  \IfFontExistsTF{Libertinus Serif}{\setmainfont{Libertinus Serif}}{}%
  \IfFontExistsTF{Libertinus Sans}{\setsansfont{Libertinus Sans}}{}%
  \@ifpackageloaded{unicode-math}{%
    \IfFontExistsTF{Libertinus Math}{\setmathfont{Libertinus Math}}{}}{}}{}
\makeatother

\usepackage{lastpage}
\usepackage{refcount}
\usepackage{tabularx}
\usepackage{tcolorbox}
\tcbuselibrary{skins,breakable}

% Folio and the continued-page note, both read from the LastPage label.
\newcommand*{\presscontinuedtext}{Continued next page}
\newcommand*{\pressfolio}{\pagemark{} / \getpagerefnumber{LastPage}}
\newcommand*{\presscontinued}{%
  \ifnum\value{page}<\getpagerefnumber{LastPage}\relax\presscontinuedtext\fi}

\makeatletter
\ifdefined\KOMAClassName
  \usepackage[automark]{scrlayer-scrpage}
  % Reserve enough room for the running header and continued footer.
  \KOMAoptions{headlines=2.2,footlines=1.3,headsepline=0.5pt,footsepline=0.5pt,plainfootsepline=true}
  \pagestyle{scrheadings}
  \clearpairofpagestyles
  \ihead{\ifdefempty{\presscode}{\headmark}{\presscode}}
  \ohead{\pressifset{\pressfields}{\pressfieldsinline}}
  \cfoot[\pressfolio]{\pressfolio}
  \ofoot[\presscontinued]{\presscontinued}
  \setkomafont{pageheadfoot}{\normalfont\footnotesize}
  \setkomafont{pagenumber}{\normalfont\footnotesize}

  % Question-level headings are easy to scan without adding colour.
  \setkomafont{disposition}{\normalfont\bfseries\color{pressink}}
  \addtokomafont{section}{\Large}
  \addtokomafont{subsection}{\large}
  \RedeclareSectionCommand[
    beforeskip=1.1\baselineskip plus 0.15\baselineskip minus 0.1\baselineskip,
    afterskip=0.45\baselineskip plus 0.1\baselineskip
  ]{section}
  \RedeclareSectionCommand[
    beforeskip=0.8\baselineskip plus 0.1\baselineskip minus 0.1\baselineskip,
    afterskip=0.35\baselineskip plus 0.05\baselineskip
  ]{subsection}
  \@ifundefined{chapter}{}{\setkomafont{chapter}{\normalfont\bfseries\LARGE\color{pressink}}}
\fi
\makeatother

\captionsetup{format=hang,skip=8pt}

\renewenvironment{quote}{%
  \def\FrameCommand{{\color{pressaccent}\vrule width 0.35mm}\hspace{0.9em}}%
  \MakeFramed{\advance\hsize-\width\FrameRestore}\noindent}
 {\endMakeFramed}

% Booktabs rules and row spacing that stay legible in print. Quarto loads
% booktabs before this file only when the document has tables.
\makeatletter
\@ifpackageloaded{booktabs}{%
  \setlength{\heavyrulewidth}{0.09em}%
  \setlength{\lightrulewidth}{0.06em}%
  \setlength{\cmidrulewidth}{0.06em}}{}
\makeatother
\AtBeginEnvironment{longtable}{%
  \renewcommand{\arraystretch}{1.12}%
  \setlength{\tabcolsep}{6pt}}

% Every callout gets the same soft grey box, with room to breathe across breaks.
\AtBeginDocument{%
  \hypersetup{linkcolor=presssecondary,urlcolor=presssecondary,citecolor=presssecondary}%
  \colorlet{quarto-callout-color}{pressaccent}%
  \colorlet{quarto-callout-color-frame}{pressaccent}%
  \colorlet{quarto-callout-note-color}{pressaccent}%
  \colorlet{quarto-callout-note-color-frame}{pressaccent}%
  \colorlet{quarto-callout-tip-color}{pressaccent}%
  \colorlet{quarto-callout-tip-color-frame}{pressaccent}%
  \colorlet{quarto-callout-warning-color}{pressaccent}%
  \colorlet{quarto-callout-warning-color-frame}{pressaccent}%
  \colorlet{quarto-callout-important-color}{pressaccent}%
  \colorlet{quarto-callout-important-color-frame}{pressaccent}%
  \colorlet{quarto-callout-caution-color}{pressaccent}%
  \colorlet{quarto-callout-caution-color-frame}{pressaccent}%
  \tcbset{
    top=4mm,
    bottom=2mm,
    right=2mm,
    boxsep=1mm,
    before skip=8pt,
    after skip=8pt,
    pad at break*=2mm,
    frame hidden,
    interior hidden,
    borderline={0.35mm}{0mm}{pressaccent},
    underlay unbroken={\path[fill=presssoft] (interior.south west) rectangle (interior.north east);},
    underlay first={\path[fill=presssoft] (interior.south west) rectangle (interior.north east);},
    underlay middle={\path[fill=presssoft] (interior.south west) rectangle (interior.north east);},
    underlay last={\path[fill=presssoft] (interior.south west) rectangle (interior.north east);}
  }}

% Write-in fields. The filter defines \pressfields as a run of
% \pressfielditem{label}; each layout below gives \pressfielditem its meaning.
\newlength{\pressfieldgap}
\setlength{\pressfieldgap}{1.6em}
\newlength{\pressfieldcell}
\newlength{\pressfieldlabel}
\newcount\pressfieldcount
\newcount\pressfieldindex

% Running head: each label followed by a short rule.
\newcommand*{\pressfieldsinline}{%
  \begingroup
  \pressfieldindex=0
  \def\pressfielditem##1{%
    \ifnum\pressfieldindex>0 \enspace\fi
    \advance\pressfieldindex by 1\relax
    ##1\enspace\underline{\hspace{3.4cm}}}%
  \pressfields
  \endgroup}

% Header block: labelled boxes sharing the text width equally.
\newcommand*{\pressfieldsboxed}{%
  \begingroup
  \pressfieldcount=0
  \def\pressfielditem##1{\advance\pressfieldcount by 1\relax}%
  \pressfields
  \setlength{\pressfieldcell}{%
    \dimexpr(\textwidth-\pressfieldgap*\numexpr\pressfieldcount-1\relax)/\pressfieldcount\relax}%
  \setlength{\fboxsep}{0.9mm}%
  \pressfieldindex=0
  \def\pressfielditem##1{%
    \ifnum\pressfieldindex>0 \hspace{\pressfieldgap}\fi
    \advance\pressfieldindex by 1\relax
    \settowidth{\pressfieldlabel}{\bfseries ##1\enspace}%
    \makebox[\pressfieldcell][l]{\bfseries ##1\enspace
      \fbox{\parbox[c][3ex][c]{\dimexpr\pressfieldcell-\pressfieldlabel-2\fboxsep-2\fboxrule\relax}{\strut}}}}%
  \noindent\pressfields\par
  \endgroup}

% The title is a ruled two-column block on the first page. Rows with nothing
% on either side are left out.
\makeatletter
\renewcommand{\maketitle}{%
  \ifdefined\KOMAClassName\thispagestyle{plain.scrheadings}\else\thispagestyle{empty}\fi
  \def\press@rows{{\Large\bfseries\presstitlefield} & \pressifset{\presslabel}{{\large\bfseries\presslabel}}\\}%
  \ifboolexpr{test {\ifdefempty{\pressseries}} and test {\ifdefempty{\pressedition}}}{}{%
    \appto\press@rows{\pressseries & \pressedition\\}}%
  \ifboolexpr{test {\ifdefempty{\presssubtitle}} and test {\ifdefempty{\@subtitle}}
      and test {\ifdefempty{\pressauthor}}}{}{%
    \appto\press@rows{\presssubtitlefield & \pressauthorfield\\}}%
  \vspace*{-\dimexpr\headheight+\headsep-0.15em\relax}%
  \noindent\rule{\textwidth}{0.6pt}\par
  \vspace{0.35em}%
  \begingroup
  \setlength{\tabcolsep}{0pt}%
  \renewcommand{\arraystretch}{1.0}%
  \noindent
  \begin{tabularx}{\textwidth}{@{}X>{\raggedleft\arraybackslash}p{0.34\textwidth}@{}}
  \press@rows
  \end{tabularx}\par
  \endgroup
  \vspace{0.7em}%
  \noindent\rule{\textwidth}{0.6pt}\par
  \pressifset{\pressfields}{\vspace{0.3em}\pressfieldsboxed}%
  \vspace{0.8em}}
\makeatother
